\section{Problem Statement}
\label{sec:problem}
% Reframed 20 Sep 2026 (author): "we are investigating quantum walks as a
% general tool". The subject is the walk as an algorithmic primitive; the
% coin is one design dimension inside it, NOT the contribution. The earlier
% framing -- "the coin operator as an unexplored degree of freedom", and the
% "design / trainability / empirical evaluation" research question -- is
% superseded; do not restore either.
%
% What this section should state:
%   * The walk is a general primitive for building quantum algorithms, and
%     universality (Childs; Lovett et al.) is what licenses treating it so.
%   * The thesis exercises that primitive across algorithm families rather
%     than developing one application: order finding, QAOA and search in
%     Ch. 3; quantum neural networks, variational classifiers and quantum
%     kernels, on graph-structured data, in Chs. 4-5.
%   * For each family the same three questions: what the walk formulation
%     is, what it costs as a circuit, what it achieves against a tuned
%     classical or quantum baseline.
%   * Design choices inside the walk -- coin (structure-dependent,
%     step-dependent, complex), shift, formulation -- are variables where an
%     application makes the comparison meaningful, not the subject.
%   * No claim of quantum advantage anywhere; simulators throughout;
%     negative results reported as results.
\subsection{Centrality of Machine Learning to Modern Computing}
\subsection{Quantum Information as a New Computational Paradigm}
\subsection{The Quantum Walk as a General Algorithmic Primitive}
% Universality, and what it does and does not license. Points forward to
% Sec. 2.4 and Ch. 3.
\subsection{What This Thesis Asks of the Primitive}
% The three questions above, the families they are asked of, and the limits
% of the claims.

